% Chapter — Mars piecewise-exponential atmosphere (FR-8, assumption A-3).
% Follows the mandatory FR-29 six-section template: purpose; assumptions
% and domain of validity; derivation; discretization and implementation;
% validation evidence; references. The equation labels eq:mars:glenn,
% eq:mars:scaleheight, and eq:mars:segment are echoed verbatim in comments
% in cpp/src/models/atmosphere_mars.cpp (FR-29 equation-label-to-code
% traceability); renaming them breaks that trace.
\chapter{Mars Atmosphere: Piecewise-Exponential Table}
\label{ch:marsatmosphere}

\section{Purpose}
\label{sec:mars:purpose}

This chapter documents the Mars atmosphere (FR-8, assumption A-3): a
piecewise-exponential density table for aerobraking-class sensitivity
studies in Mars orbit. Mars entry, descent, and landing is a project
non-goal, so the model's only downstream consumer is the cannonball drag
acceleration of Chapter~\ref{ch:drag} at orbital altitudes. The model is
implemented in \texttt{cpp/src/models/atmosphere\_mars.cpp} with its
interface in \texttt{cpp/include/star/models/atmosphere\_mars.hpp}.

\section{Assumptions and domain of validity}
\label{sec:mars:domain}

\textbf{Confidence: low --- provenance provisional per PRD A-3.} This
entire model carries the PRD's assumption-A-3 flag, prominently and
deliberately. The node densities derive from the NASA Glenn Research
Center ``Mars Atmosphere Model'' curve fits~\cite{nasaglennmars}: an
educational fit to Mars Global Surveyor data (April 1996) that states no
uncertainty, no altitude datum, no seasonal or diurnal qualification, and
whose temperature fit is discontinuous at its \qty{7}{\kilo\metre} seam.
It is a real, citable source and it is used honestly --- but it is not a
validated engineering atmosphere model. The same flag is carried in the
golden manifest (\texttt{tests/golden/atmosphere/manifest.toml}) and in
the source header. Until a citable engineering provenance replaces it
(A-3: ``reference TBD''), results that depend on absolute Mars density
are sensitivity studies, not predictions. Known context, for scale: the
real Mars atmosphere varies seasonally by tens of percent in surface
pressure and by order-of-magnitude factors in upper-atmosphere density
(dust storms); this static table represents none of that.

Assumptions:
\begin{enumerate}
  \item \textbf{Static, spherically uniform atmosphere} --- density is a
    function of altitude alone.
  \item \textbf{Piecewise-exponential structure.} Between committed
    nodes the density is exponential in altitude; the per-segment scale
    height is derived from the adjacent node densities
    (equation~\eqref{eq:mars:scaleheight}), which makes the profile
    continuous at every boundary \emph{by construction}.
  \item \textbf{Node truth equals the source curve fit sampled at
    \qty{5}{\kilo\metre} steps, 0--\qty{100}{\kilo\metre}.} The source
    temperature fit crosses absolute zero near \qty{112}{\kilo\metre},
    so the table is capped at \qty{100}{\kilo\metre}, below the
    breakdown.
\end{enumerate}

Domain of validity: altitude $z \in [\qty{0}{\kilo\metre},
\qty{100}{\kilo\metre}]$ above the source model's (unstated) surface
reference. Out-of-domain behavior, matching the code exactly: above
\qty{100}{\kilo\metre} the topmost segment's exponential is continued ---
a documented non-physical continuity aid whose scale height
($\sim\qty{41}{\kilo\metre}$, an artifact of the source fit's collapsing
temperature) is far larger than a physical Martian upper-atmosphere
scale height; it exists so that aerobraking arcs with apoapsis above the
table see a smooth, monotonically vanishing density rather than a step.
Below \qty{0}{\kilo\metre} the first segment's exponential is continued
down to \qty{-8}{\kilo\metre}, comfortably below the deepest Martian
terrain (an engineering guard, not a physical datum); below that, and
for non-finite input, the model throws \texttt{std::domain\_error}.
Both behaviors are gated by \testid{ATM-MARS-DOMAIN}.

\section{Derivation}
\label{sec:mars:derivation}

\subsection{Source curve fits and node sampling}

The NASA Glenn model~\cite{nasaglennmars}, in metric form with altitude
$h$ in metres, is
\begin{equation}
  T_C(h) =
  \begin{cases}
    -31 - 0.000998\,h, & h < 7000, \\
    -23.4 - 0.00222\,h, & h \ge 7000,
  \end{cases}
  \qquad
  p_{\mathrm{kPa}} = 0.699\,e^{-0.00009\,h},
  \qquad
  \rho = \frac{p_{\mathrm{kPa}}}{0.1921\,(T_C + 273.1)},
  \label{eq:mars:glenn}
\end{equation}
with $T_C$ in degrees Celsius and $\rho$ in
\unit{\kilogram\per\metre\cubed}. The committed nodes are
$\rho_i = \rho(z_i)$ at $z_i = 0, 5, 10, \dots, \qty{100}{\kilo\metre}$
(21 nodes), evaluated once by the golden generator and committed as exact
binary64 values; the node set, not the source fit, is the model
definition from that point on. The source's temperature discontinuity at
\qty{7}{\kilo\metre} enters only through the node values ($z = 5$ and
\qty{10}{\kilo\metre} straddle it); the committed model itself is
continuous everywhere by construction.

\subsection{Piecewise-exponential construction}

Within segment $[z_i, z_{i+1}]$ define the inverse scale length from the
adjacent nodes,
\begin{equation}
  k_i \;=\; \frac{\ln \rho_{i+1} - \ln \rho_i}{z_{i+1} - z_i}
  \qquad\left(\text{scale height } H_i = -1/k_i\right),
  \label{eq:mars:scaleheight}
\end{equation}
and evaluate
\begin{equation}
  \rho(z) \;=\; \rho_i\, e^{\,k_i\,(z - z_i)},
  \qquad z_i \le z < z_{i+1}.
  \label{eq:mars:segment}
\end{equation}
Continuity at an interior boundary $z_{i+1}$ holds by construction: the
left limit is $\rho_i \exp[k_i (z_{i+1} - z_i)] =
\rho_i \exp[\ln(\rho_{i+1}/\rho_i)] = \rho_{i+1}$ exactly in real
arithmetic, and to within a few floating-point roundings ($\sim 10^{-16}$
relative) in binary64 --- four orders below the $10^{-12}$ continuity
gate of exit criterion~9. Node exactness also holds by construction: at
$z = z_i$ the segment that starts at $z_i$ is selected, its exponent is
identically zero, $e^{0} = 1$ exactly, and the returned value is the
committed $\rho_i$ bit for bit.

\section{Discretization and implementation}
\label{sec:mars:implementation}

\begin{enumerate}
  \item \textbf{Module and traceability.} The model lives in
    \texttt{cpp/src/models/atmosphere\_mars.cpp}; the source comments
    echo the labels \texttt{eq:mars:glenn},
    \texttt{eq:mars:scaleheight}, and \texttt{eq:mars:segment} verbatim,
    and the file header carries the A-3 low-confidence flag.
  \item \textbf{Bit-exact node table.} The 21 node densities are
    compiled in as C++17 hexadecimal floating literals, numerically
    identical to the committed golden
    (\texttt{tests/golden/atmosphere/mars\_nodes.toml}, whose
    \texttt{rho\_kgpm3} hex fields are normative); test
    \testid{ATM-MARS-NODES} enforces bit equality at every node. The
    segment coefficients $k_i$ are computed once from the nodes, in
    fixed order, and cached (function-local static; single-threaded
    core, D-10).
  \item \textbf{Segment selection.} Fixed-order forward scan for the
    segment with $z_i \le z$; altitudes at or above the last node use
    the final segment's coefficient (the documented extrapolation), and
    altitudes below node~0 use segment~0's coefficient down to the
    \qty{-8}{\kilo\metre} guard.
  \item \textbf{Determinism.} Straight-line binary64 arithmetic with
    libm \texttt{log}/\texttt{exp} only; fixed evaluation order; no
    data-dependent iteration counts.
\end{enumerate}

\section{Validation evidence}
\label{sec:mars:validation}

Table~\ref{tab:mars:validation} lists the tests that constitute this
chapter's validation evidence. Each test identifier is machine-checked to
exist in the test tree by \texttt{scripts/lint\_chapters.py}; golden
provenance and tolerances are recorded in
\texttt{tests/golden/atmosphere/manifest.toml} (which repeats the A-3
low-confidence flag).

\begin{table}[htbp]
  \centering
  \caption{Validation evidence for the Mars atmosphere model (doctest,
    \texttt{cpp/tests/test\_atmosphere.cpp}).}
  \label{tab:mars:validation}
  \begin{tabular}{lp{8.6cm}}
    \toprule
    Test ID & Acceptance basis \\
    \midrule
    \testid{ATM-MARS-NODES} &
      Phase~3 exit criterion~9(a): the model evaluated at every committed
      node altitude returns the committed node density with bit-exact
      double equality (equation~\eqref{eq:mars:segment} with zero
      exponent); off-node evaluations match the independent 50-digit
      mpmath reference to $10^{-13}$ relative. \\
    \testid{ATM-MARS-CONT} &
      Phase~3 exit criterion~9(b): at every internal segment boundary,
      evaluation at the boundary and at \texttt{nextafter} increments on
      both sides yields a relative jump below $10^{-12}$. \\
    \testid{ATM-MARS-DOMAIN} &
      Out-of-domain behavior of Section~\ref{sec:mars:domain}: continued
      exponentials above \qty{100}{\kilo\metre} (monotone, vanishing)
      and down to \qty{-8}{\kilo\metre}; \texttt{std::domain\_error}
      below \qty{-8}{\kilo\metre} and for non-finite input. \\
    \bottomrule
  \end{tabular}
\end{table}

\section{References}
\label{sec:mars:references}

This chapter relies on the NASA Glenn Research Center Mars Atmosphere
Model~\cite{nasaglennmars} as the (provisional, low-confidence per PRD
A-3) source of the node densities. No other reference is load-bearing;
replacing the provenance with a validated engineering model (e.g.\ a
published Mars-GRAM-derived table) is the designated A-3 follow-up and
requires regenerating the node golden and this chapter together. Full
bibliographic details appear in the bibliography.
